\chapter{Технологическая часть}

\section{Реализация алгоритмов}
Для реализации алгоритмов использовался язык программирования C++, так как он поддерживает статическую типизацию и имеет возможность отключения сборщика мусора. Отладка происходила в среде разработки <<Visual Studio 2022>>.

% 17112

\begin{lstlisting}[style=cpp]
	void print_n_rec(int n){            // 0
		if (n == 1)                     // 1
			std::cout << n << " ";		// 2
		else{
			print_n(n - 1); 			// 3
			std::cout << n << " ";		// 4
		}
		return;                         // 5
	}
\end{lstlisting}
Листинг 1 -- Код реализации рекурсивного алгоритма с нумерацией строк.\\

\begin{lstlisting}[style=cpp]
	void print_n(int n){                // 0
		for (int i = 1; i <= n; i++)    // 1
			std::cout << i << " ";      // 2
		return;                         // 3
	}
\end{lstlisting}
Листинг 2 -- Код реализации нерекурсивного алгоритма с нумерацией строк.

\section*{Вывод}

В данном разделе приведены реализации рекурсивного и нерекурсивного алгоритмов из предыдущей работы. В этих программах отсутствуют участки, которые можно выполнить параллельно.

\clearpage
